% Ce chapitre repose sur :
%   DOC: s27, s30
%   OBS: REL-DEM.P05
%   HYP: (aucun)
%   CHI: chapitres-numerotes, modeles-total, periode-debut, periode-fin, periode-jours, profils-applicatifs
%   CHI: sections-du-rapport, source-lignes-total, source-tables, tableaux-etablissement, tdb-indicateurs, tdb-pages
%   SER: (aucune)
%   SECTIONS: 0

\chapter*{Introduction générale}
\addcontentsline{toc}{chapter}{Introduction générale}
% Un chapitre étoilé ne met pas à jour la marque de l'en-tête courant : sans cette ligne,
% les pages de ce texte portent l'en-tête de ce qui le précède — mesuré, « TABLE DES
% MATIÈRES » sur la deuxième page de l'introduction et « CHAPITRE 9 » sur celle de la
% conclusion. Le défaut ne se voit qu'en regardant le document composé.
\markboth{\MakeUppercase{Introduction générale}}{\MakeUppercase{Introduction générale}}

\section*{Contexte général}

Le stage s'est déroulé au service d'accueil et d'admission de l'hôpital provincial Sidi Saïd de
Meknès. Ce service est le point d'entrée administratif de l'hôpital. Il ouvre les dossiers,
enregistre les rendez-vous, admet et oriente, facture et recouvre. Chacun de ces gestes laisse une
trace dans le système d'information hospitalier, c'est-à-dire le logiciel par lequel
l'établissement enregistre son activité administrative et médicale.

\section*{Problématique}

Le règlement intérieur des hôpitaux range l'établissement des statistiques et la gestion de
l'information hospitalière parmi les missions propres de ce service \cite{s27}. Le responsable du
service assure le secrétariat du comité qui examine chaque mois les indicateurs d'activité. Le
comité d'établissement se prononce chaque trimestre sur le rapport d'activité, et le directeur
établit un rapport annuel. Ces trois périodicités reposent sur une seule source, les données que
le service saisit lui-même.

La liste des indicateurs attendus se lit dans la nomenclature statistique nationale \cite{s30}.
Parmi ses tableaux, \chiffre{tableaux-etablissement} descendent jusqu'à l'établissement nommé : capacité litière, journées
d'hospitalisation, admissions, taux d'occupation, durée moyenne de séjour, rotation, consultations
et activité de laboratoire.

L'écran de démarrage du logiciel porte \chiffre{profils-applicatifs} profils applicatifs, dont un profil dédié aux rapports
et aux statistiques\releve{REL-DEM.P05}. L'accès d'un stagiaire se limite au profil des
rendez-vous, et ce que produit le profil de restitution n'a pas été observé.

La question traitée est donc celle-ci. À partir des seules données que le service saisit, une
chaîne de traitement peut-elle produire les indicateurs que la réglementation attend ? Pour chaque
indicateur, il s'agit aussi de dire s'il est calculable, ce qui l'en empêche et ce qu'il vaut.

\section*{Solution envisagée}

La solution proposée est une chaîne de données complète, des fichiers extraits du logiciel jusqu'à
un tableau de bord. Elle se construit en quatre temps.

Le premier temps documente la source avant de la transformer. Chaque champ vu à l'écran est relevé
avec son identifiant, son écran et la valeur affichée.

Le deuxième construit la chaîne de traitement. Elle reçoit les fichiers bruts, type leurs valeurs,
met à part les lignes rejetées et organise le reste en tables d'analyse. Elle compte
\chiffre{modeles-total} modèles de transformation et se rejoue en entier.

Le troisième restitue les résultats dans un tableau de bord de \chiffre{tdb-pages} pages et
\chiffre{tdb-indicateurs} indicateurs. Chaque indicateur affiche sa définition, la décision qu'il
sert et les tables dont sa valeur est recalculée.

Le quatrième sépare, pour chaque résultat, ce qu'il démontre de ce qu'il établit. Un résultat
démontre que la chaîne sait produire une grandeur. Il établit le plus souvent peu de chose sur
l'établissement lui-même.

\subsection*{Les données}

Aucune donnée réelle ne sort du service, car la confidentialité des données de santé l'interdit.
La chaîne est donc bâtie sur un jeu de données synthétiques, c'est-à-dire produit par un programme
écrit pour ce projet plutôt qu'extrait du système observé. Sa volumétrie est relevée par
établissement nommé dans le recueil statistique national \cite{s30}, puis mise à l'échelle de la
période. Le jeu compte \chiffre{source-lignes-total} lignes sur \chiffre{source-tables} tables, du
\chiffre{periode-debut} au \chiffre{periode-fin}, soit \chiffre{periode-jours} jours.

Aucune grandeur produite par cette chaîne n'est donc une mesure de l'établissement. Une relation
injectée dans le jeu se retrouve nécessairement à la sortie.

\section*{Plan du rapport}

Le rapport suit l'ordre du travail, en \chiffre{chapitres-numerotes} chapitres et
\chiffre{sections-du-rapport} sections.

La première partie, \annoncePartie{Contexte et cadrage}, présente le terrain.
\annonceChapitre{1}{L'organisme d'accueil} présente l'hôpital et son service d'accueil et
d'admission. \annonceChapitre{2}{Le système d'information hospitalier} décrit le logiciel du service
et ce que le relevé des écrans a établi. \annonceChapitre{3}{Cadrage et méthodologie} expose le
besoin, la contrainte de confidentialité et la méthode du projet.

La deuxième partie, \annoncePartie{Analyse et conception}, conçoit la chaîne.
\annonceChapitre{4}{Conception du jeu de données} explique comment le jeu est calibré, quelles
relations y sont injectées et quels défauts y sont semés à dessein. \annonceChapitre{5}{Architecture
de la chaîne de données} décrit les couches traversées et l'orchestration qui les enchaîne.
\annonceChapitre{6}{Qualité des données et rapprochement d'identités} mesure la qualité des données
produites. Il évalue aussi un rapprochement probabiliste, méthode qui estime la probabilité que deux
fiches désignent la même personne.

La troisième partie, \annoncePartie{Réalisation}, restitue. \annonceChapitre{7}{Environnement et
outils} nomme les outils de la chaîne et le motif de chaque choix. \annonceChapitre{8}{Analyse de
l'activité} ventile l'activité du service et compare ce que la chaîne retrouve aux relations
injectées. \annonceChapitre{9}{Le tableau de bord} décrit l'écran de restitution, ses deux publics
et sa correspondance avec la nomenclature nationale. \annonceChapitre{10}{Recommandations} tire les
conséquences, en séparant pour chacune le constat qui la fonde de ce que l'indicateur démontre.

Une conclusion générale dit ce qui a été fait, ce qui ne l'a pas été et ce qu'il faudrait pour
porter la chaîne sur des données réelles.

Le rapport est rédigé à la forme impersonnelle, et la première personne est réservée aux passages
qui relatent le déroulement du stage.
